\Needspace{17\baselineskip}
\SourceLeaf{134}
\noindent Übrigens bemerken wir noch, dass, wenn man die Zahlen der ersten Columne als die doppelten Logarithmen von Tangenten betrachtet, die Zahlen der zweiten und dritten die doppelten Complemente der Logarithmen der dazu gehörigen Sinus und Cosinus sein werden. Man kann daher mit den gewöhnlichen trigonometrischen Tafeln den vorgesetzten Zweck ganz auf dieselbe Art erreichen, als mit den von \textsc{Leonelli} vorgeschlagenen; nur erspart theils die letztere die Division und Multiplication mit 2 (die aber doch ein mässig geübter Rechner leicht im Kopfe macht), theils gibt sie, wenn sie auf dieselbe Zahl von Decimalen berechnet ist, doppelt so viele Schärfe. Sonderbar ist es, dass die sich hierauf gründende Art, die trigonometrischen Tafeln zu gleicher Absicht anzuwenden, von \textsc{Delambre} nicht erwähnt wird, da er doch eine andere weit weniger bequeme anführt. \textsc{Leonelli} hat sehr Recht, sich zu beschweren, dass man ein solches Verfahren dem seinigen an die Seite setzen wollte.\par\medskip\begin{center}\rule{.24\linewidth}{.3pt}\end{center}

